\section{Inline Visual Route: Model Foundation}

These figures are placed in the main argument as visual proof-of-reading aids: each one separates objects, direction, quantitative or formal anchor, and evidence route.

\subsection{Carrier to realization}

The diagram below is generated from the r012 figure registry, so its objects, arrows, formula anchor, and evidence link are checked before the release audit can pass.

\begin{figure}[p]
\centering
\resizebox{0.96\textwidth}{!}{%
\begin{tikzpicture}[x=1cm,y=1cm,>=Latex,every node/.style={font=\small}]
  \definecolor{ocBlue}{HTML}{173B57}
\definecolor{ocGold}{HTML}{A77D2A}
\definecolor{ocGreen}{HTML}{4B7F52}
\definecolor{ocRed}{HTML}{8E3B32}
\definecolor{ocGray}{HTML}{ECEFF1}
  \draw[rounded corners=10pt,fill=ocGray!35,draw=ocBlue,line width=0.9pt] (-6.8,-3.4) rectangle (6.8,3.4);
  \node[font=\bfseries\large,ocBlue] at (0,3.0) {Carrier to realization};
  \node[draw,rounded corners=5pt,fill=blue!7,text width=0.24\textwidth,align=center,minimum height=1.05cm] (a) at (-4.35,1.0) {Carrier C};
  \node[draw,rounded corners=5pt,fill=green!8,text width=0.24\textwidth,align=center,minimum height=1.05cm] (b) at (4.35,1.0) {Realization R\_t};
  \draw[->,line width=1.0pt,ocGreen] (a.east) -- node[above,fill=ocGray!35,inner sep=2pt,align=center] {lawful realization} (b.west);
  \node[draw,rounded corners=5pt,fill=orange!10,text width=0.28\textwidth,align=center,minimum height=0.9cm] (r) at (-4.35,-1.45) {status at t=0,1};
  \node[draw,rounded corners=5pt,fill=white,text width=0.28\textwidth,align=center,minimum height=0.9cm] (m) at (0,-1.45) {status at t=0,1};
  \node[draw,rounded corners=5pt,fill=red!6,text width=0.28\textwidth,align=center,minimum height=0.9cm] (f) at (4.35,-1.45) {formal anchor: \( R_t=\rho(C,t) \)};
  \draw[->,dashed,ocGold] (r.north) -- (a.south);
  \draw[->,dashed,ocGold] (m.north) -- ($(a)!0.5!(b)$);
  \draw[->,dashed,ocGold] (f.north) -- (b.south);
\end{tikzpicture}}
\caption{Carrier to realization. The figure demonstrates the reader-facing route from Carrier C to Realization R\_t; the relevant variable or number is status at t=0,1, the formula anchor is R\_t=\textbackslash\{\}rho(C,t), and the evidence link is the corresponding proof, table, or replay section in the release corpus.}
\label{fig:r012-inline-28a_oc133_inline_figures_foundation-1}
\end{figure}

\subsection{Liveness status split}

The diagram below is generated from the r012 figure registry, so its objects, arrows, formula anchor, and evidence link are checked before the release audit can pass.

\begin{figure}[p]
\centering
\resizebox{0.96\textwidth}{!}{%
\begin{tikzpicture}[x=1cm,y=1cm,>=Latex,every node/.style={font=\small}]
  \definecolor{ocBlue}{HTML}{173B57}
\definecolor{ocGold}{HTML}{A77D2A}
\definecolor{ocGreen}{HTML}{4B7F52}
\definecolor{ocRed}{HTML}{8E3B32}
\definecolor{ocGray}{HTML}{ECEFF1}
  \draw[rounded corners=10pt,fill=ocGray!35,draw=ocBlue,line width=0.9pt] (-6.8,-3.4) rectangle (6.8,3.4);
  \node[font=\bfseries\large,ocBlue] at (0,3.0) {Liveness status split};
  \node[draw,rounded corners=5pt,fill=blue!7,text width=0.24\textwidth,align=center,minimum height=1.05cm] (a) at (-4.35,1.0) {Live state};
  \node[draw,rounded corners=5pt,fill=green!8,text width=0.24\textwidth,align=center,minimum height=1.05cm] (b) at (4.35,1.0) {Death boundary};
  \draw[->,line width=1.0pt,ocGreen] (a.east) -- node[above,fill=ocGray!35,inner sep=2pt,align=center] {status predicate} (b.west);
  \node[draw,rounded corners=5pt,fill=orange!10,text width=0.28\textwidth,align=center,minimum height=0.9cm] (r) at (-4.35,-1.45) {binary status};
  \node[draw,rounded corners=5pt,fill=white,text width=0.28\textwidth,align=center,minimum height=0.9cm] (m) at (0,-1.45) {binary status};
  \node[draw,rounded corners=5pt,fill=red!6,text width=0.28\textwidth,align=center,minimum height=0.9cm] (f) at (4.35,-1.45) {formal anchor: \( live(x,t)\in\{0,1\} \)};
  \draw[->,dashed,ocGold] (r.north) -- (a.south);
  \draw[->,dashed,ocGold] (m.north) -- ($(a)!0.5!(b)$);
  \draw[->,dashed,ocGold] (f.north) -- (b.south);
\end{tikzpicture}}
\caption{Liveness status split. The figure demonstrates the reader-facing route from Live state to Death boundary; the relevant variable or number is binary status, the formula anchor is live(x,t)\textbackslash\{\}in\textbackslash\{\}\{0,1\textbackslash\{\}\}, and the evidence link is the corresponding proof, table, or replay section in the release corpus.}
\label{fig:r012-inline-28a_oc133_inline_figures_foundation-2}
\end{figure}

\subsection{Residue after collapse}

The diagram below is generated from the r012 figure registry, so its objects, arrows, formula anchor, and evidence link are checked before the release audit can pass.

\begin{figure}[p]
\centering
\resizebox{0.96\textwidth}{!}{%
\begin{tikzpicture}[x=1cm,y=1cm,>=Latex,every node/.style={font=\small}]
  \definecolor{ocBlue}{HTML}{173B57}
\definecolor{ocGold}{HTML}{A77D2A}
\definecolor{ocGreen}{HTML}{4B7F52}
\definecolor{ocRed}{HTML}{8E3B32}
\definecolor{ocGray}{HTML}{ECEFF1}
  \draw[rounded corners=10pt,fill=ocGray!35,draw=ocBlue,line width=0.9pt] (-6.8,-3.4) rectangle (6.8,3.4);
  \node[font=\bfseries\large,ocBlue] at (0,3.0) {Residue after collapse};
  \node[draw,rounded corners=5pt,fill=blue!7,text width=0.24\textwidth,align=center,minimum height=1.05cm] (a) at (-4.35,1.0) {Collapse event};
  \node[draw,rounded corners=5pt,fill=green!8,text width=0.24\textwidth,align=center,minimum height=1.05cm] (b) at (4.35,1.0) {Residue class};
  \draw[->,line width=1.0pt,ocGreen] (a.east) -- node[above,fill=ocGray!35,inner sep=2pt,align=center] {classification} (b.west);
  \node[draw,rounded corners=5pt,fill=orange!10,text width=0.28\textwidth,align=center,minimum height=0.9cm] (r) at (-4.35,-1.45) {survivor count};
  \node[draw,rounded corners=5pt,fill=white,text width=0.28\textwidth,align=center,minimum height=0.9cm] (m) at (0,-1.45) {survivor count};
  \node[draw,rounded corners=5pt,fill=red!6,text width=0.28\textwidth,align=center,minimum height=0.9cm] (f) at (4.35,-1.45) {formal anchor: \( residue(x,t)>0 \)};
  \draw[->,dashed,ocGold] (r.north) -- (a.south);
  \draw[->,dashed,ocGold] (m.north) -- ($(a)!0.5!(b)$);
  \draw[->,dashed,ocGold] (f.north) -- (b.south);
\end{tikzpicture}}
\caption{Residue after collapse. The figure demonstrates the reader-facing route from Collapse event to Residue class; the relevant variable or number is survivor count, the formula anchor is residue(x,t)>0, and the evidence link is the corresponding proof, table, or replay section in the release corpus.}
\label{fig:r012-inline-28a_oc133_inline_figures_foundation-3}
\end{figure}

\subsection{Rebirth continuation}

The diagram below is generated from the r012 figure registry, so its objects, arrows, formula anchor, and evidence link are checked before the release audit can pass.

\begin{figure}[p]
\centering
\resizebox{0.96\textwidth}{!}{%
\begin{tikzpicture}[x=1cm,y=1cm,>=Latex,every node/.style={font=\small}]
  \definecolor{ocBlue}{HTML}{173B57}
\definecolor{ocGold}{HTML}{A77D2A}
\definecolor{ocGreen}{HTML}{4B7F52}
\definecolor{ocRed}{HTML}{8E3B32}
\definecolor{ocGray}{HTML}{ECEFF1}
  \draw[rounded corners=10pt,fill=ocGray!35,draw=ocBlue,line width=0.9pt] (-6.8,-3.4) rectangle (6.8,3.4);
  \node[font=\bfseries\large,ocBlue] at (0,3.0) {Rebirth continuation};
  \node[draw,rounded corners=5pt,fill=blue!7,text width=0.24\textwidth,align=center,minimum height=1.05cm] (a) at (-4.35,1.0) {Residual carrier};
  \node[draw,rounded corners=5pt,fill=green!8,text width=0.24\textwidth,align=center,minimum height=1.05cm] (b) at (4.35,1.0) {New realization};
  \draw[->,line width=1.0pt,ocGreen] (a.east) -- node[above,fill=ocGray!35,inner sep=2pt,align=center] {continuation rule} (b.west);
  \node[draw,rounded corners=5pt,fill=orange!10,text width=0.28\textwidth,align=center,minimum height=0.9cm] (r) at (-4.35,-1.45) {continuation index};
  \node[draw,rounded corners=5pt,fill=white,text width=0.28\textwidth,align=center,minimum height=0.9cm] (m) at (0,-1.45) {continuation index};
  \node[draw,rounded corners=5pt,fill=red!6,text width=0.28\textwidth,align=center,minimum height=0.9cm] (f) at (4.35,-1.45) {formal anchor: \( C_{t+1}=F(C_t) \)};
  \draw[->,dashed,ocGold] (r.north) -- (a.south);
  \draw[->,dashed,ocGold] (m.north) -- ($(a)!0.5!(b)$);
  \draw[->,dashed,ocGold] (f.north) -- (b.south);
\end{tikzpicture}}
\caption{Rebirth continuation. The figure demonstrates the reader-facing route from Residual carrier to New realization; the relevant variable or number is continuation index, the formula anchor is C\_\{t+1\}=F(C\_t), and the evidence link is the corresponding proof, table, or replay section in the release corpus.}
\label{fig:r012-inline-28a_oc133_inline_figures_foundation-4}
\end{figure}

\subsection{Type discipline}

The diagram below is generated from the r012 figure registry, so its objects, arrows, formula anchor, and evidence link are checked before the release audit can pass.

\begin{figure}[p]
\centering
\resizebox{0.96\textwidth}{!}{%
\begin{tikzpicture}[x=1cm,y=1cm,>=Latex,every node/.style={font=\small}]
  \definecolor{ocBlue}{HTML}{173B57}
\definecolor{ocGold}{HTML}{A77D2A}
\definecolor{ocGreen}{HTML}{4B7F52}
\definecolor{ocRed}{HTML}{8E3B32}
\definecolor{ocGray}{HTML}{ECEFF1}
  \draw[rounded corners=10pt,fill=ocGray!35,draw=ocBlue,line width=0.9pt] (-6.8,-3.4) rectangle (6.8,3.4);
  \node[font=\bfseries\large,ocBlue] at (0,3.0) {Type discipline};
  \node[draw,rounded corners=5pt,fill=blue!7,text width=0.24\textwidth,align=center,minimum height=1.05cm] (a) at (-4.35,1.0) {Claim type};
  \node[draw,rounded corners=5pt,fill=green!8,text width=0.24\textwidth,align=center,minimum height=1.05cm] (b) at (4.35,1.0) {Allowed evidence};
  \draw[->,line width=1.0pt,ocGreen] (a.east) -- node[above,fill=ocGray!35,inner sep=2pt,align=center] {admissibility check} (b.west);
  \node[draw,rounded corners=5pt,fill=orange!10,text width=0.28\textwidth,align=center,minimum height=0.9cm] (r) at (-4.35,-1.45) {claim family};
  \node[draw,rounded corners=5pt,fill=white,text width=0.28\textwidth,align=center,minimum height=0.9cm] (m) at (0,-1.45) {claim family};
  \node[draw,rounded corners=5pt,fill=red!6,text width=0.28\textwidth,align=center,minimum height=0.9cm] (f) at (4.35,-1.45) {formal anchor: \( claim\mapsto evidence \)};
  \draw[->,dashed,ocGold] (r.north) -- (a.south);
  \draw[->,dashed,ocGold] (m.north) -- ($(a)!0.5!(b)$);
  \draw[->,dashed,ocGold] (f.north) -- (b.south);
\end{tikzpicture}}
\caption{Type discipline. The figure demonstrates the reader-facing route from Claim type to Allowed evidence; the relevant variable or number is claim family, the formula anchor is claim\textbackslash\{\}mapsto evidence, and the evidence link is the corresponding proof, table, or replay section in the release corpus.}
\label{fig:r012-inline-28a_oc133_inline_figures_foundation-5}
\end{figure}

\subsection{Boundary test}

The diagram below is generated from the r012 figure registry, so its objects, arrows, formula anchor, and evidence link are checked before the release audit can pass.

\begin{figure}[p]
\centering
\resizebox{0.96\textwidth}{!}{%
\begin{tikzpicture}[x=1cm,y=1cm,>=Latex,every node/.style={font=\small}]
  \definecolor{ocBlue}{HTML}{173B57}
\definecolor{ocGold}{HTML}{A77D2A}
\definecolor{ocGreen}{HTML}{4B7F52}
\definecolor{ocRed}{HTML}{8E3B32}
\definecolor{ocGray}{HTML}{ECEFF1}
  \draw[rounded corners=10pt,fill=ocGray!35,draw=ocBlue,line width=0.9pt] (-6.8,-3.4) rectangle (6.8,3.4);
  \node[font=\bfseries\large,ocBlue] at (0,3.0) {Boundary test};
  \node[draw,rounded corners=5pt,fill=blue!7,text width=0.24\textwidth,align=center,minimum height=1.05cm] (a) at (-4.35,1.0) {Assumption};
  \node[draw,rounded corners=5pt,fill=green!8,text width=0.24\textwidth,align=center,minimum height=1.05cm] (b) at (4.35,1.0) {Falsifier};
  \draw[->,line width=1.0pt,ocGreen] (a.east) -- node[above,fill=ocGray!35,inner sep=2pt,align=center] {stress test} (b.west);
  \node[draw,rounded corners=5pt,fill=orange!10,text width=0.28\textwidth,align=center,minimum height=0.9cm] (r) at (-4.35,-1.45) {negative control};
  \node[draw,rounded corners=5pt,fill=white,text width=0.28\textwidth,align=center,minimum height=0.9cm] (m) at (0,-1.45) {negative control};
  \node[draw,rounded corners=5pt,fill=red!6,text width=0.28\textwidth,align=center,minimum height=0.9cm] (f) at (4.35,-1.45) {formal anchor: \( \partial\Omega \neq \varnothing \)};
  \draw[->,dashed,ocGold] (r.north) -- (a.south);
  \draw[->,dashed,ocGold] (m.north) -- ($(a)!0.5!(b)$);
  \draw[->,dashed,ocGold] (f.north) -- (b.south);
\end{tikzpicture}}
\caption{Boundary test. The figure demonstrates the reader-facing route from Assumption to Falsifier; the relevant variable or number is negative control, the formula anchor is \textbackslash\{\}partial\textbackslash\{\}Omega \textbackslash\{\}neq \textbackslash\{\}varnothing, and the evidence link is the corresponding proof, table, or replay section in the release corpus.}
\label{fig:r012-inline-28a_oc133_inline_figures_foundation-6}
\end{figure}
